\section{Conclusions}
\label{iv:conclusiones}
\begingroup
\setlength{\baselineskip}{12pt}
\setlength{\parskip}{0pt}
\setlength{\abovedisplayskip}{5pt}
\setlength{\belowdisplayskip}{5pt}

The discrete structure of the continuum developed in this article
supports an exceptional realisation reaching \(V^\natural\) and
an operator-theoretic realisation preserving compact duality.
The connection is constructed: the APP--TRIT--TPK state transmits
orientation, sheets, residues, quotients, boundary data and memory
to its numerical coordinates, incidence register, and phase and action
representations. The central contribution relates these realisations
through codes, gluings, vertex products, projectors and intertwining
operators with specified domains.

The joint construction of the continuum supplies the antecedent of
these maps. The same extensions produce the measure on histories,
the sheetwise Gram identity with its terminal remainder, the cancellation
defect, rectangular incidence and the natural transport of the memory
complex. The assembly retains all continuations and satisfies \(D^2=0\)
on its generators; its determinant line retains degrees and bases.
Compatible weights, isometric transport, the survival remainder and
the conditions for a scalar determinant reading retain the specific
roles assigned to them by the proofs. These correlated operations
constitute the common structure preceding numerical, exceptional
and dimensional representations.

% BEGIN R27_FRAMING_CONCLUSION_EN
The arithmetic and incidence connection is specified by results with
explicit domains. The defect laws distinguish modulus refinement,
periodic repetition, and increasing dimension. The aggregate operator
preserves an integral sublattice of index two. In its quadratic block
\(3H\), the normalized unit \(H\) has norm one; the Pell recurrence
determines \(H^n\), while the original block evolves as \(3^nH^n\).
The bijection between pairs
of trits and nonadic digits preserves carry, while projection of the
ambient cube has nine elements per fibre. The alphabet of \(42\) blocks,
regional junctions, and correlations of ordered records make explicit
the information preceding linear representations.

% BEGIN R28_FRAMING_CONCLUSION_EN
The five regional closures attain the minimum of eight in the finite
problem of covering the four prescribed edges. The graph retains
emission labels, and the quotient retains their incidences between
phase classes. Oriented reversal makes the reader \(\nu_{120}\) odd
and preserves the even parity of \(\nu_{270}\); these identities
make explicit the response of memory to a change of orientation.
% END R28_FRAMING_CONCLUSION_EN

% BEGIN S27_FRAMING
The regional comparison determines three polarised configurations:
the stars of \(B_K\), \(B_\varphi\) and \(B_{\rm APP}\) have
signatures \((3,1,0)\), \((1,0,3)\) and \((1,1,2)\), respectively.
The corona \(\{i:K_i\geq729\}=B_K=H_e\cap P_\pi\) links
the numerical register to an incidence four-subset, and the signed
cochain fixes \(H^-_\alpha=B_K\sqcup\{5,7\}\).
The complete census places the signature \((3,1,0)\) at fifteen
four-subsets. Rigidity belongs to the ordered frame
\((H^-_\alpha,H_e,H_\varphi,H_{\rm APP})\): its stabiliser
in \(G_W\) is the identity. This result fixes the permutation
freedom on the twelve positions; the values of \(K\) retain
their determination by the integer register and its reconstruction
operations.

The integer lift of the regional signature supplies a complementary
result. The identities \(p+e+f=q+3c\) and \(p+e-f=a+3h\) retain
the quotients accompanying the residues, and the carry cocycle makes
their composition compatible with association of the sum.
In the complete table of \(27\) columns, the triples \((0,0,0)\)
and \((1,2,0)\) have \(q=a=0\), whereas their pairs \((c,h)\)
are \((0,0)\) and \((1,1)\). Each reduction \(\bar c,\bar h\),
adjoined separately to the evaluation of \(1,p,e,f\), increases
the rank from four to five. The complete reading \((q,a,c,h)\)
itself retains fibres: \((1,2,0)\) and \((2,1,0)\) produce the
same datum \((0,0,1,1)\). The transport identities and cylinder
bounds therefore specify both the additional information retained
by integer memory and the multiplicity of antecedents remaining
in each reading. This articulation jointly preserves the numerical
content and incidence of the regional history.

% END S27_FRAMING


The identity \(C^2=5I_6\) follows from quadratic correlation evaluated
at the five finite positions. Saturation and dual neutrality leave a
pair of boundaries whose representative depends on the declared
orientation. Closure controls distinguish that selection from variations
of the boundary and of register position. 

The exceptional result follows the incidence of \(K\) through to the
vertex algebra. The dodecaphasic register is constructed before its
rational reading and support selection. Its signed configuration enters
the Witt flag; the ternary code determines the gluing of \(A_2^{12}\),
and the flag with the specified metric, origin and chamber determines
the Leech neighbour. The order-two and order-three extensions of its
lattice algebra realise \(V^\natural\) by application of the
corresponding construction theorems, whose hypotheses are verified
in the lattice corridor. The chosen isomorphisms transport the vacuum,
conformal vector, grading, product and traces; they do not identify an
order-two automorphism with an order-three automorphism. The Monster
acts on \(V^\natural\), and Moonshine determines its graded traces.
The rational reading of the register and the digits of the constants
have other domains: their connection with the algebra is the constructed
chain of objects and maps, not an action on their digits.

In the exceptional realization,
the rational identity proves \(J-(t_3+12)\in3^9q\mathbb Z[[q]]\);
its first coefficient determines optimality of the uniform exponent.
The formal proof retains its scope at every degree, while truncated
checks retain their role as finite comparisons.
% END R27_FRAMING_CONCLUSION_EN

The same \(K\) determines a geometric polarisation and a recovery
reference. In the fixed real representation, the nonzero vector
\(u_K\) determines a line in \(V_{11}\); its exact squared norm
normalizes the rank-one term of the rank-ten orthogonal projector,
whose covariance law is specified. The full
dodecaphasic orbit is invertible and identifies operators from their
responses with exact stability bounds. Its rational encoding preserves
the twelve blocks in the fixed chart; bilateral balance recovers the
state from two channels defined by isometries, and its unitary extension jointly
recovers state and incoming memory. Invariance of other observables
retains the proved commutation condition. The arithmetic--geometric
role of \(K\) therefore comprises incidence, direction and identification
in addition to its scalar evaluation.

The unital extension of the Weyl representation establishes a continuous
realisation of tracial dimension. The phase and shift operators generate
each finite matrix algebra, and full-fibre inclusions preserve the unit
and normalised trace. The nonadic uniformly hyperfinite closure has a
tracial representation whose weak closure is a hyperfinite factor of
type \(II_1\); normalised ranks attain every tracial dimension in the
unit interval within it. A marked continuation uses another inclusion
and produces a different closure. The result identifies the refinement
operation responsible for preserving global normalisation.

The action and angular pair determine the compact form rather than
receiving a prescribed radius. Normalised regional renewal and the
positive action sections fix the common scale of the ellipse; its
angular coordinates fix the semiaxis ratio. The unique positive radius
satisfies
\[
 R_\star^4=\frac{A-C^*}{A+C^*}
 =\frac{499\alpha-\eta_{\mathrm{ret}}}{501\alpha+\eta_{\mathrm{ret}}}.
\]
Membership in the elliptic chart and the choice of positive root are
proved. The semiaxis product preserves action; their exchange transforms
the radius into \(R_\star^{-1}\). The angular charts, axis orientation
and action scale thus retain their respective roles in the quadratic
form.

\Needspace{12\baselineskip}
The established T-duality is a scalar and operator equivalence.
A single involution preserves the compact form and unitarily conjugates
the self-adjoint Hamiltonians of reciprocal radii, including their
domains and quadratic-form domains. Covariance of the residual section
requires transport of \(L_9\); on the corresponding quotient, the
minimum energy is well defined and preserves duality. String
normalisation transports these identities to dimensional radii without
changing the generated shape radius. The elliptic collection contains both
reciprocal fibres and preserves the restriction of domains and of the
screen isomorphism. Modular semiconjugacy on the imaginary axis connects
their parameters and functions while retaining the proper domain of
the operators in each realisation.

The screens express the dimensional hierarchy through specified maps.
The reductions \(12\to11\to10\) and \(26\to24\) arise, respectively,
from removal of the uniform mode and the direction \(u_K\), and from
orthogonal restriction followed by the isotropic quotient of the
Lorentzian lattice. The joint morphism identifies the quotient
presentations and intertwines the compact sector. The graph retains
the coordinate in \(V_{11}\) together with its projection; each
individual projection retains its kernel. Ternary puncturing adds
the length-eleven combinatorial screen with distance five and a perfect
partition. These operations determine the ranks, rather than an
identification between spaces of coincident dimension.

The extension towards M-theory has a determined algebraic support
and a precise dynamical condition. Changes of direction and arithmetic
type along trajectories, composed with their preceding boundary datum,
form an exactly additive cocycle. The worldsheet, Polyakov action and
codomain of eleven-dimensional fields subsequently receive this
structure through the realisation maps. The transport theorem preserves
\(Q^2=H\) on the image of the isometry under the stated intertwining
and domain conditions. Complete physical identification includes the
fields, tensions and action normalisation, the oscillator sector and
its amplitudes, anomaly control and the low-energy gravitational limit
of each compactification
(Section~\ref{iv:condiciones-realizacion-fisica}).
These requirements fix the dynamics into which the proved structures
are transported.

A common chain is thus established linking the discrete construction
of the continuum with the Moonshine algebra, dodecaphasic polarisation,
tracial realisation, compact duality and action composition.
Its maps determine what information remains, which quotient is taken
and which equations are preserved. This is the mathematical content
received by the dimensional realisations and on which the developments
of statistics, radiation and species structure continue.

\par\endgroup
